% SOURCE_AUTHORITY: sga5_fr_workpass.tex, lines 15266--15274 (LF-normalized unit).
% SOURCE_SHA256: 791F4EFFC5E02832D5D77ED03518C8156D6F07E4C8238B03545DB93D883FBB28
Demonstração do lema 2: como $M=T(L)$ é projetivo em cada grau, a hipótese do lema é expressa dizendo que $T(\varphi)$ é homotópico a $\varphi_0$; existe, portanto, um operador de homotopia $k_0=(k_0^i)_i$, com $k_0^i:M^i\to M^{i-1}$, tal que
\[
T(\varphi)-\varphi_0=d_0k_0+k_0d_0,
\]
onde $d_0$ é a diferencial de $M$. Pelo lema 3(i), $k_0^i$ levanta-se a uma aplicação linear $k^i:L^i\to L^{i-1}$. Obtém-se assim um operador de homotopia $k=(k^i)_i$ em $L$. Basta pôr
\[
\varphi'=\varphi-kd-dk
\]
($d=$ diferencial de $L$) para verificar o lema 2.
